\chapter{Strategic and Emerging Risk Assessment}
\label{chap:Strategic and Emerging Risk Assessment}

%---------------------------- SECTION ----------------------------
\section{The Risks That Change Everything}
\paragraph{Every risk examined in the preceding chapters of Part VI — health and safety, environmental, financial, cyber, third-party, reputational, people and conduct — operates, in the main, within a relatively stable understanding of what the organisation is, what it does, and the environment in which it operates. These are risks to the execution of the strategy — risks that, if they materialise, disrupt, damage, or constrain the organisation's ability to pursue its objectives within a broadly understood world.}
\paragraph{\textbf{Strategic risk} is different in character. It is the risk that the strategy itself is wrong — that the organisation's fundamental assumptions about its markets, its competitive position, its business model, its regulatory environment, or its relevance are mistaken, incomplete, or becoming obsolete. Strategic risk does not operate within the existing framework — it threatens the framework itself. And when it materialises, the consequences are not operational disruption or regulatory sanction but something more fundamental: the erosion of the organisation's reason for being, the destruction of its competitive position, or, in the most serious cases, the loss of its viability as a going concern.}
\paragraph{\textbf{Emerging risk} is equally distinctive in character. It is the risk that is not yet fully formed — that is visible, perhaps, as a distant signal on the horizon, or detectable only through the aggregation of weak signals that individually appear insignificant but collectively point towards something important. Emerging risks have not yet manifested in observable harm. They may not be well understood — the mechanisms through which they could cause harm may still be uncertain, the timeline may be unclear, and the probability may be genuinely unknown rather than merely difficult to estimate. But they are real — and the organisations that identify and engage with them early are dramatically better placed than those that wait until they are obvious.}
\paragraph{For compliance and legal professionals, strategic and emerging risk assessment is perhaps the area where the profession's contribution is most underutilised and most significant. Compliance functions that focus exclusively on the operational and regulatory dimensions of risk — important as those are — miss the opportunity to contribute to the strategic conversations that determine the organisation's long-term direction and resilience. A compliance professional who can engage credibly with strategic risk — who can identify the regulatory and legal dimensions of strategic choices, who can surface the emerging regulatory trends that will shape the strategic environment, and who can help the organisation think clearly about the risks it is taking at the level of strategy — is a more valuable partner to leadership than one whose contribution is limited to the current compliance framework.}
\paragraph{This chapter explores what strategic and emerging risk assessment involves, how it differs from other forms of risk assessment, and how to build the capability to do it well.}

%---------------------------- SECTION ----------------------------
\section{What Is Strategic Risk?}
\paragraph{The concept of strategic risk is used in different ways in different organisations and different frameworks — and it is worth being precise about what it means before exploring how to assess it.}
\paragraph{At its most fundamental, strategic risk is the risk of significant loss of value arising from strategic decisions — from the choices the organisation makes about what to do, how to do it, in what markets to operate, and on what business model to build. It encompasses:}
\paragraph{\textbf{Strategy execution risk} — the risk that the organisation fails to execute its chosen strategy effectively — through inadequate capabilities, insufficient resources, poor implementation, or organisational resistance. This is the most familiar form of strategic risk and the one most commonly addressed in strategic planning processes.}
\paragraph{\textbf{Strategy choice risk} — the risk that the chosen strategy is itself wrong — that the assumptions on which it is built are flawed, that it misreads the competitive environment, or that it fails to anticipate the strategic moves of competitors or the evolution of customer needs. This is a harder and more uncomfortable form of strategic risk to assess — because it requires the organisation to challenge its own strategic direction rather than simply managing the risks of executing it.}
\paragraph{\textbf{Business model risk} — the risk that the organisation's fundamental business model — the way it creates, delivers, and captures value — becomes obsolete, disrupted, or unviable. Business model risk is often a slow-building strategic risk — the signs are visible long before the crisis, but the organisational and psychological investment in the existing model makes it difficult to respond until the disruption is well advanced.}
\paragraph{\textbf{Strategic concentration risk} — the risk of over-dependence on a single market, a single product, a single customer, a single channel, or a single geography — creating vulnerability to adverse developments in that single dimension that the organisation's diversification does not provide resilience against.}
\paragraph{\textbf{Competitive and market risk} — the risk that competitive dynamics — new entrants, disruptive technologies, changing customer preferences, or shifts in the regulatory framework — erode the organisation's competitive position faster than it can adapt}
\paragraph{\textbf{Regulatory and political strategic risk} — the risk that changes in the regulatory or political environment fundamentally alter the conditions under which the organisation operates — rendering existing business models unviable, creating new barriers to entry in current markets, or opening new opportunities that the organisation is not positioned to exploit.}
\paragraph{\textbf{Macro-environmental risk} — the risk that broad macroeconomic, demographic, technological, or geopolitical shifts create conditions that are materially adverse for the organisation's strategy — whether through economic recession, demographic change, technological disruption, or geopolitical instability.}

%---------------------------- SECTION ----------------------------
\section{Why Strategic Risk Is Different}
\paragraph{Strategic risk assessment is qualitatively different from the operational, financial, and compliance risk assessment that forms the bulk of most organisations' risk management activity — and understanding those differences is essential for designing an approach that is genuinely effective.}

\section{The Time Horizon Challenge}
\paragraph{Most operational risk assessment operates over a relatively short time horizon — the coming year, the next regulatory cycle, the next planning period. Strategic risk operates over much longer horizons — the shifts in competitive dynamics, regulatory environments, and technological landscapes that create strategic risk typically unfold over years or decades rather than months.}
\paragraph{This time horizon creates a specific challenge for risk assessment. Risks that are genuinely significant over a five or ten-year horizon may appear to be manageable or distant when assessed against the near-term indicators that dominate most risk monitoring frameworks. The organisations that have been most seriously damaged by strategic risk — Kodak by digital photography, Blockbuster by streaming, Nokia by the smartphone — had early warning signals available to them long before the disruption was irreversible. The challenge was not the absence of information — it was the inability to act on weak signals that were inconsistent with the existing strategic narrative.}

\section{The Uncertainty Challenge}
\paragraph{Most risk assessment — even qualitative risk assessment — works within a framework where the nature of the risk is broadly understood, where historical data provides some basis for likelihood estimation, and where the range of potential consequences can be described with reasonable confidence. Strategic risk assessment frequently operates in conditions of genuine uncertainty — where not just the probability but the fundamental nature of future outcomes is unclear.}
\paragraph{This distinction between \textbf{risk} — where probabilities can be assigned to a defined set of outcomes — and \textbf{uncertainty} — where the outcome space itself cannot be fully defined — is the economist Frank Knight's famous distinction, and it matters enormously for assessment methodology. Quantitative risk assessment techniques that work well in conditions of risk perform poorly or misleadingly in conditions of genuine uncertainty. The appropriate analytical response to uncertainty is not more precise modelling but more robust scenario thinking — accepting that the future cannot be predicted and focusing instead on developing genuine resilience across a range of possible futures.}

\section{The Cognitive and Political Challenge}
\paragraph{Strategic risk assessment faces a specific set of cognitive and political challenges that are absent or less prominent in operational risk assessment. These challenges arise from the fact that strategic risk assessment necessarily involves challenging the organisation's own strategic direction — and the people most responsible for that direction are often the most prominent voices in strategic risk discussions.}
\paragraph{\textbf{Confirmation bias} — the tendency to seek evidence that confirms existing beliefs — is particularly powerful in the strategic context. Management teams that have committed publicly to a strategy, that have invested significant resources in its execution, and whose careers and reputations are bound up with its success are not naturally disposed to give adequate weight to evidence that the strategy is wrong.}
\paragraph{\textbf{The innovator's dilemma} — Clayton Christensen's concept of the conditions under which successful organisations fail to respond to disruptive innovation — describes precisely the dynamics of strategic risk blindness. Successful organisations, serving their best customers with their best products, are structurally resistant to the low-margin, low-quality innovations that eventually displace them — because those innovations look, at first, like inferior alternatives that only appeal to less attractive customers.}
\paragraph{\textbf{Horizon fixation} — the organisational tendency to focus attention on the risks that are most immediately visible and most clearly quantifiable, at the expense of the more distant, more uncertain, and more qualitative risks that may be more strategically significant. A risk management framework that is highly effective at identifying and managing operational and financial risks can actually compound strategic risk blindness — by creating a sense of security that the risks that matter are being managed, whilst the strategic environment shifts around it.}

%---------------------------- SECTION ----------------------------
\section{What Is Emerging Risk?}
\paragraph{Emerging risk — sometimes called horizon risk — is risk that is new, developing, or insufficiently understood to be fully incorporated into existing risk assessments. Emerging risks have several distinctive characteristics:}
\paragraph{\textbf{They are visible before they are quantifiable} — emerging risks can often be detected as weak signals — early indicators of potential future harm — long before they have developed sufficiently to be assessed with any precision. The ability to identify and engage with weak signals before they become strong ones is the essence of emerging risk management.}
\paragraph{\textbf{They develop through interconnected systems} — many of the most significant emerging risks arise not from a single new hazard but from the interaction of multiple developments — technological, social, regulatory, and environmental — that individually appear manageable but collectively create novel exposures. The emerging risk of AI-enabled fraud, for example, arises from the intersection of advances in artificial intelligence, the increasing digitisation of financial transactions, and the growing sophistication of criminal organisations — no single one of which is itself a new risk.}
\paragraph{\textbf{They evolve rapidly} — by definition, emerging risks are in the process of developing. Their nature, their probability, and their potential consequences may change significantly over relatively short periods — requiring assessment frameworks that can update quickly rather than relying on periodic review cycles.}
\paragraph{\textbf{They may have no historical precedent} — for truly novel emerging risks — those arising from genuinely new technologies, genuinely new social dynamics, or genuinely unprecedented combinations of conditions — the historical data that conventional risk assessment relies on may simply not exist. Assessment in these conditions requires qualitative judgement, scenario thinking, and analogical reasoning rather than quantitative modelling.}

%---------------------------- SECTION ----------------------------
\section{The Strategic and Emerging Risk Assessment Process}
\paragraph{Applying the core risk assessment process to strategic and emerging risk requires some significant adaptations — both in the identification techniques used and in the evaluation and management approaches adopted.}

\section{Identification — Scanning the Horizon}
\paragraph{The identification of strategic and emerging risks requires a different approach from the process-based hazard identification appropriate for operational risks. It requires \textbf{systematic environmental scanning} — a structured, ongoing process of monitoring the organisation's external environment for developments that could create new risks or fundamentally alter existing ones.}
\paragraph{\textbf{PESTLE analysis} — the systematic consideration of Political, Economic, Social, Technological, Legal, and Environmental factors — is one of the most widely used frameworks for strategic environmental scanning. Used rigorously and regularly, PESTLE analysis provides a structured prompt for identifying developments across the full range of the external environment that could affect the organisation's strategic position.}
\paragraph{Applied to risk assessment, a PESTLE analysis should ask:}
\paragraph{\textbf{Political} — what political developments — domestically and internationally — could affect the organisation's operating environment? Changes of government, shifts in regulatory philosophy, geopolitical tensions, trade policy changes, and political instability in key markets can all create strategic risks with significant compliance dimensions.}
\paragraph{\textbf{Economic} — what macroeconomic conditions and trends could affect the organisation's financial performance and strategic viability? Interest rate trajectories, inflation dynamics, economic growth prospects, exchange rate trends, commodity price cycles, and the financial health of key markets and sectors are all relevant.}
\paragraph{\textbf{Social} — what changes in social values, demographic patterns, consumer behaviour, and societal expectations could affect demand for the organisation's products and services, or create new expectations about how it should conduct itself? The growing expectations around ESG performance, the shifting demographic profile of customers and employees, and the evolving expectations around digital experience are all social trends with strategic risk implications.}
\paragraph{\textbf{Technological} — what technological developments could disrupt the organisation's business model, create new competitive threats, change customer expectations, or alter the regulatory landscape? Artificial intelligence, quantum computing, blockchain, biotechnology, and the accelerating pace of digital transformation are all technological developments with significant strategic risk dimensions for organisations across many sectors.}
\paragraph{\textbf{Legal} — what changes in the legal and regulatory environment — domestic, international, and emerging — could alter the conditions under which the organisation operates? New legislation, regulatory reform, evolving judicial interpretations, and the development of international standards can all create strategic legal risks that go beyond the immediate compliance implications.}
\paragraph{\textbf{Environmental} — what environmental and climate-related developments could affect the organisation's strategic position? Physical climate risks, transition risks, biodiversity-related regulatory changes, and the evolving expectations of investors and customers around environmental performance are all strategic environmental risk dimensions explored in Chapter~\ref{chap:Environmental and Sustainability Risk Assessment}.}

\section{The VUCA Framework — Navigating Volatile Environments}
\paragraph{The concept of \textbf{VUCA} — developed originally in a military context and widely adopted in strategic management — provides a useful framework for characterising the conditions under which strategic risk assessment is conducted.}
\paragraph{\textbf{Volatility} — the pace and unpredictability of change in the environment. High volatility means that conditions are changing rapidly and in ways that are difficult to predict — requiring risk frameworks that are adaptive and responsive rather than static.}
\paragraph{\textbf{Uncertainty} — the degree to which the future is knowable. High uncertainty means that even well-informed analysis cannot reliably predict outcomes — requiring approaches that build resilience across scenarios rather than optimising for a predicted future.}
\paragraph{\textbf{Complexity} — the number and interconnectedness of the factors that shape outcomes. High complexity means that simple cause-and-effect analysis is insufficient — requiring systems thinking that addresses the interactions between factors rather than analysing them in isolation.}
\paragraph{\textbf{Ambiguity} — the degree to which the meaning and implications of observed developments are unclear. High ambiguity means that the same information can be interpreted in multiple ways — requiring processes that actively seek diverse perspectives and that challenge confident interpretations.}
\paragraph{For compliance professionals engaged in strategic risk assessment, VUCA provides a useful diagnostic: where is the organisation's strategic environment most volatile, most uncertain, most complex, or most ambiguous? Those are the areas where conventional risk assessment is most likely to fail and where the most careful, most nuanced approach is most needed.}

\section{Scenario Planning — The Primary Tool for Strategic Risk}
\paragraph{If the risk matrix is the primary tool for operational risk assessment, scenario planning is the primary tool for strategic risk assessment. We have encountered scenario analysis several times in this book — as a tool for evaluating tail risks in operational contexts, as the mandated approach for climate-related financial risk, and as a general analytical technique in Chapter~\ref{chap:Risk Assessment Tools and Methodologies}. In the strategic risk context, it takes on its fullest and most powerful expression.}
\paragraph{\textbf{Strategic scenario planning} involves constructing a small number of distinct, plausible, and internally consistent narratives about how the future could unfold — and examining the implications of each scenario for the organisation's strategy, its business model, and its risk profile. The purpose is not to predict which scenario will occur — it is to understand how the organisation would be affected by different futures, and to develop strategies that are robust across the range of plausible scenarios rather than optimised for a single predicted outcome.}
\paragraph{Effective scenario planning involves several stages:}
\paragraph{\textbf{Identifying the key uncertainties} — the factors about the future that are both most uncertain and most consequential for the organisation's strategic position. These might include the pace of technological disruption in the organisation's sector, the direction of regulatory reform, the evolution of customer behaviour, or the trajectory of geopolitical relationships. The key uncertainties are the variables around which the scenario narratives are constructed.}
\paragraph{\textbf{Constructing the scenarios} — developing a small number — typically two to four — of distinct scenarios that represent meaningfully different futures, constructed around different combinations of how the key uncertainties resolve. The scenarios should be plausible — grounded in real trends and reasonable extrapolations — and they should be genuinely distinct — representing different enough futures that they create different strategic challenges and require different strategic responses.}
\paragraph{\textbf{Stress testing the strategy} — examining how the organisation's current strategy performs under each scenario. Where does the strategy perform well across all scenarios — indicating robust strategic choices? Where does it perform well under some scenarios but poorly under others — indicating strategic choices that are well-suited to some futures but vulnerable to others? Where does it perform poorly under most scenarios — indicating fundamental strategic vulnerabilities?}
\paragraph{\textbf{Identifying strategic options} — developing strategic choices that improve the organisation's performance across scenarios — either by hedging against scenario-specific vulnerabilities, by building capabilities that are valuable across multiple futures, or by maintaining the strategic flexibility to adapt as the future becomes clearer.}
\paragraph{\textbf{Building early warning indicators} — identifying the observable signals that would indicate which scenario is unfolding — so that the organisation can monitor the environment and update its strategic response as the picture becomes clearer}
\paragraph{For compliance and legal professionals, scenario planning is particularly valuable in several contexts:}
\begin{itemize}
    \bitem{Regulatory scenario planning}{examining how the regulatory environment could evolve under different political and policy scenarios, and assessing the strategic implications for the organisation's compliance obligations and business model.}
    \bitem{Technology disruption scenarios}{exploring how technological developments — AI, blockchain, quantum computing — could alter the competitive landscape, the regulatory framework, and the nature of the compliance challenge.}
    \bitem{Geopolitical scenarios}{assessing the strategic and compliance implications of different geopolitical outcomes — trade tensions, sanctions regimes, regulatory divergence across jurisdictions.}
    \bitem{Climate transition scenarios}{as discussed in Chapter~\ref{chap:Environmental and Sustainability Risk Assessment}, examining how different climate pathway scenarios affect the organisation's strategic position, its asset values, and its regulatory obligations.}
\end{itemize}

\section{Weak Signal Detection — Finding Emerging Risks Early}
\paragraph{The identification of truly emerging risks — those that have not yet developed sufficiently to be visible in conventional risk assessment processes — requires a specific capability: the ability to detect and interpret \textbf{weak signals} — early, fragmentary, or ambiguous indicators of potential future risks.}
\paragraph{Weak signals are, by their nature, easy to miss and easy to dismiss. They appear as isolated data points rather than clear trends. They may be inconsistent with the prevailing strategic narrative. They may be visible only at the edges of the organisation's usual information environment — in conversations with unusual customers, in developments in adjacent industries, or in the academic and specialist literature that most busy managers do not have time to read.}
\paragraph{Building the capability to detect and act on weak signals requires:}
\paragraph{\textbf{Broad information sourcing} — actively seeking out information from a wide range of sources — including those outside the organisation's usual information diet. Academic research, regulatory consultations in adjacent sectors, technology conference proceedings, NGO reports, investigative journalism, and conversations with unusual or non-mainstream stakeholders can all provide weak signal intelligence that does not appear in standard management information.}
\paragraph{\textbf{Diverse perspectives} — weak signals that are invisible from within the organisation's mainstream perspective are often visible from the periphery — from people who interact with the organisation's environment from different vantage points. Front-line staff, unusual customers, external advisers, and board members with diverse backgrounds can all see things that the core management team misses.}
\paragraph{\textbf{Psychological safety for unconventional thinking} — surfacing weak signals requires people to be willing to raise observations that may seem speculative, inconsistent with the prevailing view, or even threatening to established ways of doing things. The conditions of psychological safety discussed in Chapter~\ref{chap:Building a Risk-Aware Culture} are as important for emerging risk identification as they are for operational risk management.}
\paragraph{\textbf{Structured reflection time} — weak signal detection does not happen in the midst of operational activity. It requires deliberate, protected time for reflection — for reading broadly, for engaging with external perspectives, and for thinking about what observed signals might mean for the organisation's future. Organisations that do not create space for this kind of reflective thinking are systematically blind to the signals that are most important for their strategic future.}
\paragraph{\textbf{Dedicated horizon scanning processes} — some organisations — particularly in fast-moving sectors — establish formal horizon scanning functions or processes that are specifically tasked with monitoring the external environment for emerging risks and reporting their findings to governance bodies. Where this is done well — with genuine intellectual rigour, genuine commitment to challenging the prevailing view, and genuine board engagement — it can provide significant early warning value.}

\section{The Delphi Method — Harvesting Expert Judgement}
\paragraph{For emerging risks where expert knowledge is widely distributed and where no single individual has a complete picture, the \textbf{Delphi Method} — a structured process for aggregating expert judgements through iterative anonymous questioning — can be a valuable assessment tool.}
\paragraph{In a Delphi process, a panel of experts is asked independently to estimate the likelihood, timing, and significance of a defined emerging risk. Their responses are aggregated and fed back to the panel anonymously — allowing each expert to see the range of views without knowing who holds them. The experts are then asked to revise their assessments in light of the group's aggregated responses. This iterative process continues until the responses converge or the range of disagreement is clearly documented.}
\paragraph{The Delphi Method is particularly valuable for emerging risks because it:}
\begin{itemize}
  \item{Aggregates expertise from multiple domains — which is often necessary where an emerging risk arises from the interaction of multiple fields.}
  \item{Reduces the influence of dominant personalities and status hierarchies on the assessment — which are particularly significant in strategic risk discussions.}
  \item{Makes explicit the range of expert disagreement — providing a more honest representation of genuine uncertainty than a single point estimate would suggest.}
\end{itemize}

\section{Red Teaming — Challenging the Strategic Narrative}
\paragraph{\textbf{Red teaming} — borrowed from military intelligence and increasingly adopted in corporate strategy and risk contexts — involves assigning a group explicitly to challenge and attempt to disprove the organisation's strategic assumptions and risk assessments. The red team's mandate is to find the flaws in the prevailing view — to identify the strategic risks that the organisation's leadership may be too close, too invested, or too cognitively biased to see.}
\paragraph{Red teaming is particularly valuable for strategic risk assessment because it directly addresses the confirmation bias and groupthink dynamics that make strategic risk blindness so common. A well-conducted red team exercise can surface strategic vulnerabilities that no amount of conventional strategic planning would reveal — precisely because its mandate is to challenge rather than confirm.}
\paragraph{For compliance and legal professionals, the red team role — challenging the prevailing view, asking the uncomfortable questions, and maintaining the independence of mind needed to see risks that others prefer not to acknowledge — is in many ways a description of the compliance function's core strategic contribution. The compliance professional who is genuinely independent — who can challenge strategic assumptions, identify the regulatory risks in commercially attractive proposals, and maintain analytical integrity in the face of commercial pressure — is performing an invaluable red team function for the organisation.}

%---------------------------- SECTION ----------------------------
\section{Regulatory and Legal Dimensions of Strategic Risk}
\paragraph{For compliance and legal professionals, the regulatory and legal dimensions of strategic risk deserve specific attention — both because they are a significant source of strategic risk in their own right, and because understanding them is the most direct contribution the compliance function can make to strategic risk assessment.}

\section{Regulatory Strategic Risk}
\paragraph{\textbf{Regulatory strategic risk} — the risk that changes in the regulatory environment fundamentally alter the strategic landscape — is one of the most significant strategic risks for organisations in regulated sectors. The history of financial services, healthcare, energy, and other heavily regulated sectors is rich with examples of regulatory changes that transformed competitive dynamics, rendered existing business models unviable, or created entirely new strategic opportunities for firms that were well positioned.}
\paragraph{The compliance function's horizon scanning activity — described in Chapter~\ref{chap:Step 1 — Identifying Hazards and Risk Sources} — is the primary vehicle through which regulatory strategic risk is identified. But translating regulatory intelligence into strategic risk assessment requires a capability that goes beyond the identification of specific new requirements — it requires the ability to understand the strategic implications of regulatory change for the organisation's business model, competitive position, and long-term viability.}
\paragraph{Key regulatory strategic risk areas that are particularly prominent in the current environment include:}
\paragraph{\textbf{The digitalisation agenda} — regulatory frameworks around digital assets, open banking, embedded finance, and digital identity are reshaping the competitive landscape across financial services in ways that create both strategic risk and strategic opportunity. Firms that are positioned to benefit from regulatory reform in these areas — and firms whose existing advantages may be eroded — need to assess these developments as strategic risks, not just compliance challenges}
\paragraph{\textbf{The sustainability regulatory agenda} — as discussed in Chapter~\ref{chap:Environmental and Sustainability Risk Assessment}, the rapidly developing body of sustainability regulation — CSRD, CSDDD, TCFD, and the emerging frameworks around nature and biodiversity — is creating strategic risk for organisations whose business models are most exposed to transition risk, and strategic opportunity for those best positioned to meet the new expectations of regulators, investors, and customers.}
\paragraph{\textbf{The AI regulatory agenda} — the EU AI Act, the developing UK approach to AI governance, and the emerging international frameworks around AI safety and accountability are creating a regulatory landscape for artificial intelligence that will significantly affect the strategic viability of AI-dependent business models. Understanding the regulatory trajectory — and its implications for the organisation's AI strategy — is a direct compliance function contribution to strategic risk assessment.}
\paragraph{\textbf{Data and privacy regulation} — the continuing development of data protection regulation globally — the GDPR's ongoing evolution, the proliferation of national data protection frameworks, and the emerging frameworks around data localisation and data sovereignty — creates strategic risks for data-intensive business models that depend on the ability to process and transfer personal data across jurisdictions.}
\paragraph{\textbf{Competition and market structure} — regulatory attention to market power, algorithmic pricing, data monopolies, and the competitive dynamics of platform markets is creating strategic uncertainty for organisations whose competitive position depends on network effects, data advantages, or market dominance that may attract regulatory challenge.}

\section{Geopolitical and Sanctions Risk}
\paragraph{\textbf{Geopolitical risk} — the risk that political developments in the international environment create conditions that are materially adverse for the organisation's strategy — has increased significantly in prominence over the past decade and is now one of the most significant sources of strategic uncertainty for internationally active organisations}
\paragraph{The compliance dimension of geopolitical risk is substantial — through the rapidly evolving \textbf{sanctions regime} that is increasingly a primary instrument of geopolitical competition. The breadth and complexity of the sanctions landscape — US OFAC sanctions, EU sanctions, UK sanctions, and the complex interaction between them — has increased dramatically, and the legal and regulatory risk of sanctions non-compliance has become one of the most significant financial crime compliance risks for any internationally active organisation.}
\paragraph{Beyond sanctions, geopolitical risk creates strategic risk through:}
\begin{itemize}
    \bitem{Trade policy uncertainty}{tariffs, trade restrictions, export controls, and the fragmentation of global supply chains create strategic risk for organisations with internationally integrated operations.}
    \bitem{Regulatory divergence}{the divergence of regulatory frameworks between jurisdictions — most visibly between the EU, the UK, and the US — creates strategic complexity for organisations that operate across multiple jurisdictions.}
    \bitem{Geopolitical conflict and instability}{armed conflict, political instability, and the risk of state failure in key markets create direct operational and strategic risks that require assessment in the enterprise risk framework.}
\end{itemize}

%---------------------------- SECTION ----------------------------
\section{Emerging Risk — Key Areas for the Current Environment}
\paragraph{Having established the framework for emerging risk assessment, it is worth examining several of the most significant emerging risk areas in the current environment — both to illustrate the assessment approach and to equip compliance professionals with an awareness of the risks most likely to shape their professional landscape in the years ahead.}

\section{Artificial Intelligence Risk — Beyond the Current Framework}
\paragraph{The regulatory framework for AI — introduced in Chapter~\ref{chap:Technology and Data in Risk Assessment} — represents the regulatory system's current understanding of AI risk. But the pace of AI development means that the risks of tomorrow's AI systems may be substantially different from — and in some cases substantially more significant than — those addressed by today's regulatory frameworks.}
\paragraph{\textbf{Large language model risks} — the rapid development of large language models and generative AI systems creates emerging risks that are still being understood: the risks of AI-generated misinformation at scale, the risks of AI-enabled social engineering and fraud, the risks of over-reliance on AI systems in high-stakes decision-making, and the systemic risks of widespread dependence on a small number of foundational AI models whose failure or compromise could create cascading consequences.}
\paragraph{\textbf{Autonomous decision-making risks} — as AI systems take on more consequential decision-making roles — in credit, in healthcare, in criminal justice, in national security — the risks of algorithmic error, algorithmic bias, and the erosion of meaningful human oversight become increasingly significant. The compliance implications of these developments — for explainability, accountability, and the protection of individual rights — are still being worked through.}
\paragraph{\textbf{AI and the compliance function itself} — as AI tools are increasingly adopted within compliance functions — for transaction monitoring, regulatory change management, document review, and risk assessment support — the governance of those tools, and the risks of over-reliance on them, become compliance risks in their own right.}

\section{Quantum Computing Risk}
\paragraph{\textbf{Quantum computing} — still in its development phase but advancing at a pace that is beginning to generate genuine near-term concern — poses an emerging risk to the cryptographic foundations of modern cybersecurity. Many of the encryption algorithms that currently protect financial transactions, personal data, and sensitive communications are vulnerable to attack by sufficiently powerful quantum computers.}
\paragraph{The timeline for this risk remains uncertain — estimates of when quantum computers will be capable of breaking current encryption vary widely. But the implications of the risk — and the length of time required to migrate to quantum-resistant cryptographic standards — mean that organisations need to begin assessing and planning for this transition now. The \textbf{National Cyber Security Centre (NCSC)} and equivalent bodies in other jurisdictions are already providing guidance on quantum-safe cryptography — a clear indicator that this emerging risk is approaching the point of requiring active management.}

\section{Climate Tipping Points and Physical Risk Acceleration}
\paragraph{As discussed in Chapter~\ref{chap:Environmental and Sustainability Risk Assessment}, the physical risks of climate change are typically modelled as gradual, linear processes — incremental increases in temperature, sea level, and extreme weather frequency. But climate science increasingly identifies the risk of \textbf{tipping points} — thresholds beyond which climate dynamics change qualitatively rather than incrementally, potentially triggering rapid and irreversible shifts that go far beyond the scenarios that current physical risk models capture.}
\paragraph{The potential activation of climate tipping points — the collapse of the West Antarctic ice sheet, the dieback of the Amazon rainforest, the disruption of major ocean circulation systems — represents an emerging risk of a genuinely novel character: one where the historical data and physical models that currently underpin climate risk assessment may be significantly underestimating the potential severity of future impacts.}
\paragraph{For compliance professionals with responsibilities around climate risk assessment and disclosure, the emerging science around climate tipping points is material — and ignoring it in favour of models based exclusively on smooth transition scenarios may represent a significant understatement of physical risk.}

\section{Pandemic and Biological Risk}
\paragraph{The COVID-19 pandemic — as a risk event — was not entirely unforeseeable. Pandemic risk had been identified and assessed by many organisations, and the scenario was well known in public health and emergency management communities. What was widely underestimated was the speed, the scale, and the systemic nature of the disruption — and the degree to which the pandemic would expose vulnerabilities in global supply chains, financial systems, and organisational operating models that had not been adequately stress tested.}
\paragraph{The emerging risk dimension of pandemic and biological risk includes:}
\begin{itemize}
    \item{The possibility of future pandemic events with different characteristics — higher mortality rates, different transmission dynamics, or greater antibiotic or antiviral resistance.}
    \item{The emerging risk of engineered biological threats — whether through accident or deliberate action — as advances in biotechnology make the modification of biological agents more accessible.}
    \item{The systemic vulnerabilities exposed by the COVID-19 pandemic that have not yet been fully addressed — in supply chain resilience, in business continuity planning, and in the governance of organisations whose operating models are highly dependent on physical presence and global mobility.}
\end{itemize}

\section{Demographic and Social Structural Risk}
\paragraph{Long-term demographic and social structural changes — whilst not typically described as emerging risks — are creating strategic risks for organisations across many sectors that are only beginning to be fully appreciated.}
\paragraph{\textbf{Demographic ageing} in many developed economies — the combination of falling birth rates and rising life expectancy that is producing a significant shift in the age profile of populations — creates strategic risks for pension systems, healthcare organisations, retail businesses, and labour markets that are still in the process of unfolding. Organisations whose business models, workforce strategies, and product portfolios are calibrated for the demographic profile of the past may find them increasingly misaligned with the demographic profile of the future.}
\paragraph{\textbf{Generational value shifts} — the different values, expectations, and priorities of younger generations relative to their predecessors — around work, around consumption, around financial services, and around organisational purpose — are creating strategic risks for organisations that have not adapted their products, their cultures, and their conduct standards to the changing expectations of their most important future stakeholders.}
\paragraph{\textbf{Social cohesion and inequality risk} — growing evidence of social cohesion challenges in many societies — rising inequality, political polarisation, declining trust in institutions — creates a strategic risk environment that is more volatile, more unpredictable, and more susceptible to sudden discontinuities than a more stable social environment would be.}

%---------------------------- SECTION ----------------------------
\section{The Compliance Function's Contribution to Strategic Risk Assessment}
\paragraph{Having mapped the landscape of strategic and emerging risk, it is worth bringing this chapter — and Part VI as a whole — to a close by reflecting on the compliance and legal function's specific contribution to strategic risk assessment.}
\paragraph{The compliance function is not the primary owner of strategic risk — that responsibility sits with the board and executive leadership. But the compliance function has a set of capabilities and perspectives that make it a uniquely valuable contributor to strategic risk conversations — and compliance professionals who develop and exercise those capabilities are significantly more valuable to their organisations than those who confine themselves to the operational and regulatory dimensions of their role}

\section{Regulatory Intelligence as Strategic Intelligence}
\paragraph{The compliance function's horizon scanning capability — its systematic monitoring of regulatory developments, enforcement trends, and supervisory priorities — produces intelligence that is as strategically relevant as it is operationally important. A regulatory change that will fundamentally alter the competitive landscape, an enforcement trend that signals a tightening of conduct expectations, or an emerging international regulatory framework that will create new compliance obligations — all of these are simultaneously compliance intelligence and strategic intelligence.}
\paragraph{Ensuring that regulatory intelligence flows into the organisation's strategic risk assessment process — not just into the operational compliance programme — is one of the most valuable contributions the compliance function can make to strategic governance.}

\section{Legal Risk as Strategic Signal}
\paragraph{The legal function's visibility of litigation trends, regulatory investigations, contractual disputes, and legal developments across the organisation's operating environment provides a rich source of strategic signal intelligence. The emergence of a new theory of liability, the success of a regulatory enforcement action in a new area, or a pattern of litigation affecting peer organisations are all potential indicators of emerging strategic risk — and ensuring that legal intelligence informs strategic risk assessment is a natural extension of the legal function's core role.}

\section{The Challenge Function — Maintaining Strategic Integrity}
\paragraph{Perhaps the most important strategic contribution of the compliance and legal function is the one that is most difficult to exercise and most easy to neglect: the willingness and the capability to challenge strategic decisions and assumptions when they carry unacceptable risk.}
\paragraph{The compliance professional who points out, clearly and persistently, that a commercially attractive strategy carries regulatory risks that have not been adequately assessed — who challenges the assumptions behind a business case, who identifies the emerging regulatory signals that contradict the strategic narrative, or who insists that a novel product or service be properly stress-tested against regulatory expectations before launch — is fulfilling the most valuable strategic risk function that the second line can provide.}
\paragraph{This requires courage — the courage to deliver messages that are unwelcome, to maintain positions that are unpopular, and to prioritise long-term risk management over short-term commercial convenience. It also requires credibility — the credibility that comes from consistently high-quality analysis, genuine understanding of the business, and a track record of adding value rather than simply imposing constraints.}
\paragraph{The compliance and legal professional who has earned that credibility — and who exercises that courage — is one of the most important strategic assets an organisation has.}

%---------------------------- SECTION ----------------------------
\newpage
\section{Chapter Summary}
In this chapter we learned about the following key points:

\begin{table}[H]
  \centering
  \renewcommand{\arraystretch}{1.5}
  \begin{tabularx}{\textwidth}{ | >{\raggedright\arraybackslash}p{0.20\textwidth} 
								| >{\raggedright\arraybackslash}p{0.40\textwidth} 
                                | >{\raggedright\arraybackslash}X | }
    \hline
    \textbf{Risk Type} & \textbf{Key Characteristics} & \textbf{Primary Assessment Approach} \\
    \hline
    \textbf{Strategy execution risk} & Failure to deliver the chosen strategy & Strategic project risk assessment; capability gap analysis\\
    \hline
    \textbf{Strategy choice risk} & The chosen strategy is fundamentally wrong & Scenario planning; red teaming; external challenge\\
    \hline
    \textbf{Business model risk} & Core business model becomes obsolete or disrupted & Disruption scenario analysis; competitive intelligence; horizon scanning\\
    \hline
    \textbf{Regulatory strategic risk} & Regulatory change alters the strategic landscape & Regulatory horizon scanning; impact assessment; scenario analysis\\
    \hline
    \textbf{Geopolitical risk} & Political developments create adverse strategic conditions & Geopolitical scenario analysis; sanctions risk assessment\\
    \hline
    \textbf{Emerging risk} & Novel, poorly understood, or not yet fully formed & Weak signal detection; Delphi method; PESTLE; VUCA analysis\\
    \hline
  \end{tabularx}
  \caption{Risk Types}
\end{table}

\begin{table}[H]
  \centering
  \renewcommand{\arraystretch}{1.5}
  \begin{tabularx}{\textwidth}{ | >{\raggedright\arraybackslash}p{0.35\textwidth} 
								| >{\raggedright\arraybackslash}p{0.35\textwidth} 
                                | >{\raggedright\arraybackslash}X | }
    \hline
    \textbf{Emerging Risk} & \textbf{Key Compliance Dimensions} & \textbf{Assessment Priority} \\
    \hline
    \textbf{AI and large language models} & Explainability; bias; governance; fraud enablement & High and rising\\
    \hline
    \textbf{Quantum computing} & Cryptographic vulnerability; data security & Medium now; rising rapidly\\
    \hline
    \textbf{Climate tipping points} & Physical risk underestimation; disclosure adequacy & High\\
    \hline
    \textbf{Pandemic and biological risk} & Business continuity; supply chain resilience & Medium\\
    \hline
    \textbf{Demographic and social structural change} & Consumer protection; conduct standards; workforce & Medium to long term\\
    \hline
  \end{tabularx}
  \caption{Emerging Risks}
\end{table}

\begin{table}[H]
  \centering
  \renewcommand{\arraystretch}{1.5}
  \begin{tabularx}{\textwidth}{ | >{\raggedright\arraybackslash}p{0.20\textwidth} 
								| >{\raggedright\arraybackslash}p{0.40\textwidth} 
                                | >{\raggedright\arraybackslash}X | }
    \hline
    \textbf{Tool} & \textbf{Primary Value} & \textbf{When Most Useful} \\
    \hline
    \textbf{PESTLE analysis} & Structured environmental scanning & Periodic strategic review; horizon scanning\\
    \hline
    \textbf{Scenario planning} & Exploring a range of plausible futures & Major strategic decisions; regulatory change assessment\\
    \hline
    \textbf{Weak signal detection} & Identifying emerging risks before they are obvious & Continuous; horizon scanning function\\
    \hline
    \textbf{Red teaming} & Challenging strategic assumptions and narratives & Pre-decision challenge; annual strategic review\\
    \hline
    \textbf{Delphi method} & Aggregating expert judgement under uncertainty & Novel or highly uncertain emerging risks\\
    \hline
    \textbf{VUCA framework} & Characterising the strategic risk environment & Calibrating the overall approach to strategic risk\\
    \hline
  \end{tabularx}
  \caption{Strategic Risk Assessment Tools}
\end{table}

\paragraph{Key takeaways:}
\begin{itemize}
  \item{Strategic risk operates at the level of the organisation's fundamental assumptions and business model — not just the execution of its strategy — and requires assessment approaches that can challenge those assumptions rather than simply managing the risks within them.}
  \item{The time horizon of strategic risk assessment is fundamentally longer than that of operational risk assessment — and the weak signals of future strategic risks are often visible long before the risks themselves are obvious.}
  \item{Scenario planning is the primary tool for strategic risk assessment — its value lies not in predicting the future but in building strategies that are robust across a range of plausible futures.}
  \item{Emerging risks require specific identification and assessment capabilities — including weak signal detection, diverse information sourcing, psychological safety for unconventional thinking, and the ability to act on uncertainty rather than waiting for certainty.}
  \item{The compliance function's most strategic contributions are regulatory intelligence — translating regulatory developments into strategic risk intelligence — and the challenge function — maintaining the independence and the courage to challenge strategic assumptions when they carry unacceptable risk.}
  \item{The organisations most damaged by strategic risk are rarely those that lacked the information — they are those that could not bring themselves to act on it. Building the governance, the culture, and the professional capability to act on weak signals and uncomfortable insights is the most important investment in strategic risk management.}
\end{itemize}

\barquote{With Chapter 24 complete, Part VI — Applied Risk Assessment — is finished. We have now covered the full landscape of the book: from the conceptual foundations of Part I, through the regulatory framework of Parts II and III, through the complete five-step process of Part III, through communication and culture in Part IV, through tools and technology in Part V, and through the full range of applied risk assessment contexts in Part VI.}

\barquote{In Part VII — the final section of the book — we bring everything together: exploring how to build and sustain an integrated risk assessment programme, how to demonstrate its value to stakeholders, and what the future of risk assessment looks like for the compliance and legal professionals who will shape it..}